\chapter{Solution Algorithms}
\label{chap:algorithms}

This chapter describes the four POMDP solvers used in this thesis and explains why they were chosen.
Section~\ref{sec:scope} defines the scope of the solver set.
Section~\ref{sec:exact} describes the exact solvers accessed through the \texttt{pomdp-solve} C library.
Section~\ref{sec:approx} describes the approximate solvers implemented in Python.
Section~\ref{sec:justify} justifies the choice of these solvers for the Foggy Forest benchmark.
Section~\ref{sec:summary3} summarizes.


\section{Scope and Solver Set}
\label{sec:scope}

This thesis studies four solvers: Witness, Incremental Pruning (IP), QMDP, and Point-Based Value Iteration (PBVI).
They represent two distinct approaches to POMDP solving.
Witness and IP are exact algorithms that compute the optimal value function up to convergence.
QMDP and PBVI are approximate algorithms that sacrifice optimality for tractability.

The four solvers cover a useful range.
Exact solvers show what the optimal solution looks like on small problems.
Approximate solvers show what is achievable on larger ones.
Comparing the two groups on the same benchmark reveals the cost of approximation and the benefit of scalability.

The solver set was chosen to be minimal but informative.
Adding more solvers would complicate the comparison without adding much to the analysis.
These four solvers represent the main families that appear in the POMDP literature \cite{shaniSurveyPointbasedPOMDP2013}.


\section{Exact Solvers via \texttt{pomdp-solve}}
\label{sec:exact}

The exact solvers are accessed through the \texttt{pomdp-solve} library, a well-established C implementation of several exact POMDP algorithms.
Cassandra et al. \cite{cassandra1994acting} developed the original solver and have maintained it over many years.
Using this library rather than reimplementing the algorithms avoids introducing bugs in the core solver logic.

The Witness algorithm was one of the first practical exact POMDP solvers.
It performs exact dynamic programming by finding the alpha vectors that are needed to represent the optimal value function at each step.
It uses a ``witness point'' procedure to enumerate these vectors efficiently.

Incremental Pruning (IP) is another exact algorithm that improves on earlier approaches by interleaving the pruning of dominated vectors with the enumeration step \cite{incrementalPruningNotes}.
This reduces the number of intermediate vectors that must be kept in memory.
Both Witness and IP produce the same optimal solution in the end; they differ in computational efficiency.

The C implementation is mature and well-tested.
Interfacing with it from Python through a subprocess or wrapper is straightforward.
This approach is also commonly used in the POMDP research community \cite{gaoTechnicalQASite2021}.


\section{Approximate Solvers in Python}
\label{sec:approx}

The two approximate solvers, QMDP and PBVI, are implemented directly in Python.
Python was chosen because it offers easy integration with numerical libraries and visualization tools.
Both algorithms are conceptually simple enough that a clean Python implementation is feasible without significant performance loss for the problem sizes considered here.

QMDP \cite{zhangModelApproximationScheme1997} is a fast heuristic that approximates the POMDP value function by assuming the current uncertainty will resolve after the next step.
It solves the underlying MDP first and then uses the resulting value function to choose actions.
The policy is:
\begin{equation}
  \pi_\text{QMDP}(b) = \arg\max_{a \in A} \sum_{s \in S} Q^*(s, a)\, b(s),
\end{equation}
where $Q^*(s, a)$ is the optimal MDP Q-value.
QMDP runs very fast but can fail when information gathering is important, because it does not model the value of reducing uncertainty.

PBVI \cite{pineauPointbasedValueIteration} maintains a set of belief points and performs exact backups only at those points.
The algorithm starts with a small set of beliefs and adds more over time.
For each belief $b$ in the set, the update computes:
\begin{equation}
  \alpha^*(b) = \arg\max_{\alpha \in \Gamma'} \alpha \cdot b,
\end{equation}
where $\Gamma'$ is the set of candidate vectors produced by the backup.
By restricting updates to sampled beliefs, PBVI avoids the exponential blowup of the full alpha vector set \cite{spaanPointbasedPOMDPAlgorithm2004}.
The resulting policy is near-optimal over the sampled region of the belief space.


\section{Why These Solvers for Foggy Forest}
\label{sec:justify}

The Foggy Forest environment is a grid-based navigation benchmark with partial observability.
The agent moves on a 2D grid and receives noisy observations about its surroundings.
The environment is small enough for exact solvers to work, but large enough to make the computational cost of exact methods visible.

Using exact solvers on this benchmark lets us measure the growth of the alpha vector set as the horizon increases.
This directly illustrates the computational challenge described in Chapter~2.
The exact solution also serves as a reference point for evaluating the quality of the approximate solutions.

QMDP and PBVI were chosen because they represent two different approximation strategies.
QMDP is a policy heuristic: it ignores future uncertainty and plans as if the state were observable.
PBVI is a value function approximation: it covers the belief space with sample points and computes local backups.
Comparing these two on Foggy Forest reveals how each approximation affects solution quality and convergence.

The combination of C-based exact solvers and Python-based approximate solvers reflects a practical engineering trade-off.
Exact solvers are fast because they are written in C, which matters when the alpha vector set is large.
Approximate solvers are implemented in Python because they need to integrate tightly with the belief update logic, visualization module, and experiment pipeline.
This division of responsibility keeps the codebase manageable while still supporting meaningful comparisons.


\section{Summary}
\label{sec:summary3}

This chapter described the four solvers used in this thesis.
Witness and IP provide exact solutions and are accessed through the \texttt{pomdp-solve} C library.
QMDP and PBVI provide approximate solutions and are implemented in Python.

The choice of solvers reflects the goal of the thesis: to study POMDP solver behavior on a controlled benchmark.
The exact solvers establish what the optimal solution looks like, and the approximate solvers show what is achievable at lower computational cost.
Chapter~4 describes how these solvers are integrated into the analysis tool and applied to the Foggy Forest environment.
